\chapter{Hasil dan Pembahasan}

\section{Hasil Implementasi atau Pelaksanaan Penelitian}
\TemplateTodo{Laporkan apa yang benar-benar berhasil dibuat, dikumpulkan, atau dijalankan. Pisahkan hasil faktual dari interpretasi.}

% Pseudocode ditempatkan di Bab IV, bukan Bab III, karena yang dilaporkan adalah
% prosedur yang benar-benar dijalankan. Karena itu pula naskah proposal tidak
% memuat algorithm2e sama sekali.
%
% \KwInput, \KwOutput, dan \KwReturn adalah versi dwibahasa yang didefinisikan
% config/thesis-preamble.tex: keduanya mencetak Masukan/Keluaran/Kembalikan pada
% naskah Indonesia dan Input/Output/Return pada naskah Inggris. Perintah bawaan
% algorithm2e (\KwIn, \KwOut, \Return) tetap tersedia, tetapi selalu berbahasa
% Inggris. Nomornya menjadi "Algoritma 4. 1" dan masuk ke DAFTAR ALGORITMA.
\begin{algorithm}[H]
\caption{Contoh format algoritma penelitian}\label{alg:contoh}
\KwInput{Data atau masukan penelitian}
\KwOutput{Hasil yang dievaluasi}
Lakukan tahap persiapan data\;
Terapkan metode yang diusulkan\;
\ForEach{skenario pengujian}{
  Evaluasi keluaran menggunakan metrik yang sesuai\;
}
\KwReturn{Ringkasan hasil dan analisis}
\end{algorithm}

Prosedur tersebut diringkas pada Algoritma~\ref{alg:contoh}.

\section{Hasil Pengujian}
\TemplateTodo{Sajikan hasil utama sesuai skenario pada Bab III. Tabel harus memiliki judul jelas, satuan, dan keterangan yang cukup.}

Hasil utama dirangkum pada Tabel \ref{tab:hasil-pengujian}.

% CONTOH TABEL 3: tabel angka, kolom sempit, tanpa kolom lentur.
%
% Ketika seluruh kolom berisi angka pendek, tabularx tidak diperlukan: cukup
% longtable dengan kolom l dan c biasa, yang lebarnya menyesuaikan isi. Kolom c
% dipakai untuk angka agar rata tengah dan mudah dibandingkan antarbaris.
%
% \multicolumn{4}{|l|}{\textbf{...}} membuat baris judul kelompok yang
% membentang di seluruh lebar tabel, dipakai untuk memisahkan blok skenario
% tanpa harus memecah tabel menjadi beberapa tabel kecil.
\begin{longtable}{|l|c|c|c|}
\caption{Contoh tabel hasil pengujian.}\label{tab:hasil-pengujian}\\
\hline
Skenario & Metrik 1 & Metrik 2 & Metrik 3 \\
\hline
\endfirsthead
\hline
Skenario & Metrik 1 & Metrik 2 & Metrik 3 \\
\hline
\endhead
\hline
\endlastfoot
\multicolumn{4}{|l|}{\textbf{Protokol pengujian utama}} \\
\hline
Skenario A & 0,00 & 0,00 & 0,00 \\
\hline
Skenario B & 0,00 & 0,00 & 0,00 \\
\hline
\multicolumn{4}{|l|}{\textbf{Protokol ablasi}} \\
\hline
Skenario C & 0,00 & 0,00 & 0,00 \\
\end{longtable}

Angka desimal ditulis dengan koma pada naskah Bahasa Indonesia. Di dalam mode matematika, tulis \verb|$0{,}00$| agar tidak muncul spasi setelah komanya.

% CONTOH TABEL 4: tabel lebar dengan banyak kolom.
%
% Tabel yang kolomnya banyak akan meluber melewati margin. Bungkus dengan
% \begingroup ... \endgroup lalu perkecil ukuran huruf dan jarak antarkolom di
% dalamnya. Karena dibungkus grup, kedua pengaturan itu tidak bocor ke tabel
% berikutnya.
%
% \cline{3-6} memberi garis hanya pada kolom 3 sampai 6, dipakai ketika beberapa
% baris berbagi isi kolom pertama sehingga kolom itu dibiarkan menyatu.
\begingroup\footnotesize\setlength{\tabcolsep}{4pt}
\begin{xltabular}{\linewidth}{|P{0.06\linewidth}|P{0.16\linewidth}|P{0.20\linewidth}|Y|P{0.08\linewidth}|P{0.08\linewidth}|}
\caption{Contoh tabel lebar dengan ukuran huruf diperkecil.}\label{tab:hasil-lebar}\\
\hline
Kode & Model & Faktor & Nilai & MAE & RMSE \\
\hline
\endfirsthead
\hline
Kode & Model & Faktor & Nilai & MAE & RMSE \\
\hline
\endhead
\hline
\endlastfoot
S1 & Model dasar & Konfigurasi & Nilai bawaan & 0,00 & 0,00 \\
\hline
A1a & Model usulan & Parameter yang diubah & Nilai pertama & 0,00 & 0,00 \\
\cline{4-6}
A1b & & & Nilai kedua & 0,00 & 0,00 \\
\end{xltabular}\endgroup

\section{Pembahasan}
\TemplateTodo{Jelaskan mengapa hasil tersebut terjadi, hubungannya dengan teori dan penelitian terdahulu, serta implikasinya. Jangan hanya mengulang isi tabel.}

\section{Keterbatasan Penelitian}
\TemplateTodo{Tuliskan keterbatasan secara jujur dan spesifik, misalnya keterbatasan data, waktu, perangkat, metode, validasi, atau cakupan pengguna.}
